\documentclass[11pt]{article}

\usepackage[margin=1in]{geometry}
\usepackage{amsmath,amssymb,bm}
\usepackage{tcolorbox}
\usepackage{enumitem}

\title{Displacement, Strain, and the Acoustic Wave Equation}
\author{}
\date{}

\begin{document}

\maketitle

\section{Displacement as a perturbation from equilibrium}

Consider a material particle whose equilibrium, or reference, position is

\begin{equation}
    \bm{x} = (x_1,x_2,x_3).
\end{equation}

After the medium is disturbed, the particle moves to a new position

\begin{equation}
    \bm{x}' = \bm{x} + \bm{u}(\bm{x},t),
\end{equation}

where

\begin{equation}
    \boxed{
    \bm{u}(\bm{x},t)
    =
    \bm{x}'-\bm{x}
    }
\end{equation}

is the displacement vector.

Thus, $\bm{u}$ does \emph{not} represent the absolute position of the particle with respect to the coordinate origin. Instead, it represents the perturbation of the particle away from its equilibrium position.

\begin{tcolorbox}[title=Key idea]
\[
\boxed{
\text{current position}
=
\text{equilibrium position}
+
\text{displacement}
}
\]

or, equivalently,

\[
\boxed{
\bm{x}'=\bm{x}+\bm{u}.
}
\]
\end{tcolorbox}

The coordinates $\bm{x}$ are therefore normally used to label the particle according to its reference position, while $\bm{u}(\bm{x},t)$ tells us how far that particle has moved.

\section{Why strain depends on derivatives of displacement}

The infinitesimal strain tensor is

\begin{equation}
    \boxed{
    \epsilon_{ij}
    =
    \frac{1}{2}
    \left(
    \frac{\partial u_i}{\partial x_j}
    +
    \frac{\partial u_j}{\partial x_i}
    \right).
    }
\end{equation}

This expression shows that strain does not measure displacement itself. Instead, it measures how the displacement changes from one material particle to another.

For example, suppose the entire medium undergoes a rigid translation,

\begin{equation}
    \bm{u}=\bm{c},
\end{equation}

where $\bm{c}$ is constant in space. Then

\begin{equation}
    \frac{\partial u_i}{\partial x_j}=0,
\end{equation}

and therefore

\begin{equation}
    \epsilon_{ij}=0.
\end{equation}

This is physically correct: the entire body has moved, but it has not been deformed.

\section{The one-dimensional interpretation}

In one spatial dimension,

\begin{equation}
    \epsilon_{xx}
    =
    \frac{\partial u}{\partial x}.
\end{equation}

Consider two nearby particles initially separated by a small distance $dx$. Their displacements are

\[
u(x,t)
\]

and

\[
u(x+dx,t).
\]

The difference between their displacements is approximately

\begin{equation}
    u(x+dx,t)-u(x,t)
    \approx
    \frac{\partial u}{\partial x}\,dx.
\end{equation}

Therefore:

\begin{itemize}
    \item if both particles move by the same amount, their separation does not change and the strain is zero;
    \item if their displacements are different, their separation changes and the material is compressed or extended.
\end{itemize}

Thus strain measures \emph{relative motion between neighboring particles}.

\section{Particle displacement and particle velocity}

The displacement field $\bm{u}$ also defines the particle velocity,

\begin{equation}
    \boxed{
    \bm{v}
    =
    \frac{\partial \bm{u}}{\partial t}.
    }
\end{equation}

This is the actual velocity of the material particles as the elastic or acoustic disturbance passes through the medium.

The particle acceleration is

\begin{equation}
    \frac{\partial^2\bm{u}}{\partial t^2}.
\end{equation}

These quantities should be distinguished from the wave propagation velocity, such as $V_P$ or $V_S$. The particles oscillate locally, while the wave itself propagates through the medium.

\section{Can displacement be used in a fluid?}

Yes. Although acoustic wave propagation is often described using pressure, the displacement vector $\bm{u}$ is also a perfectly valid physical variable in a fluid.

A fluid particle has an equilibrium position and can be displaced from it just like a particle in a solid:

\begin{equation}
    \bm{x}'=\bm{x}+\bm{u}(\bm{x},t).
\end{equation}

The important difference is not the existence of displacement, but the constitutive behavior of the material.

A solid can support both

\begin{itemize}
    \item compression,
    \item shear.
\end{itemize}

An ideal acoustic fluid can support compression but cannot support static shear stress.

Therefore, only the volumetric part of the deformation is important for acoustic propagation.

\section{Volumetric strain}

The divergence of the displacement,

\begin{equation}
    \nabla\cdot\bm{u},
\end{equation}

measures the local fractional change in volume for small deformations.

In Cartesian coordinates,

\begin{equation}
    \nabla\cdot\bm{u}
    =
    \frac{\partial u_x}{\partial x}
    +
    \frac{\partial u_y}{\partial y}
    +
    \frac{\partial u_z}{\partial z}.
\end{equation}

If

\begin{equation}
    \nabla\cdot\bm{u}>0,
\end{equation}

the local volume has increased.

If

\begin{equation}
    \nabla\cdot\bm{u}<0,
\end{equation}

the medium has been compressed.

For an acoustic fluid, the pressure perturbation is related to the volumetric strain by

\begin{equation}
    \boxed{
    p=-K\,\nabla\cdot\bm{u},
    }
\end{equation}

where $K$ is the bulk modulus.

The minus sign has a simple physical meaning: compression corresponds to

\[
\nabla\cdot\bm{u}<0,
\]

and therefore produces a positive pressure perturbation,

\[
p>0.
\]

\section{Equation of motion in a fluid}

Newton's second law for an acoustic fluid can be written as

\begin{equation}
    \boxed{
    \rho
    \frac{\partial^2\bm{u}}{\partial t^2}
    =
    -\nabla p,
    }
\end{equation}

where $\rho$ is density.

Substituting

\begin{equation}
    p=-K\nabla\cdot\bm{u}
\end{equation}

gives

\begin{equation}
    \rho
    \frac{\partial^2\bm{u}}{\partial t^2}
    =
    \nabla\left(K\nabla\cdot\bm{u}\right).
\end{equation}

Thus the acoustic wave equation can be written directly in terms of particle displacement:

\begin{equation}
    \boxed{
    \rho
    \frac{\partial^2\bm{u}}{\partial t^2}
    =
    \nabla\left(K\nabla\cdot\bm{u}\right).
    }
\end{equation}

For a homogeneous medium, $K$ and $\rho$ are constants. Defining

\begin{equation}
    c^2=\frac{K}{\rho},
\end{equation}

we obtain

\begin{equation}
    \frac{\partial^2\bm{u}}{\partial t^2}
    =
    c^2
    \nabla\left(\nabla\cdot\bm{u}\right).
\end{equation}

\section{Why acoustic displacement is longitudinal}

An acoustic fluid does not support shear waves. Therefore, acoustic particle motion is longitudinal: the particle displacement is locally associated with compression and expansion rather than shear deformation.

For a purely acoustic displacement field,

\begin{equation}
    \nabla\times\bm{u}=0.
\end{equation}

Using the vector identity

\begin{equation}
    \nabla^2\bm{u}
    =
    \nabla(\nabla\cdot\bm{u})
    -
    \nabla\times(\nabla\times\bm{u}),
\end{equation}

the condition

\begin{equation}
    \nabla\times\bm{u}=0
\end{equation}

gives

\begin{equation}
    \nabla(\nabla\cdot\bm{u})
    =
    \nabla^2\bm{u}.
\end{equation}

Consequently, in a homogeneous acoustic medium,

\begin{equation}
    \boxed{
    \frac{\partial^2\bm{u}}{\partial t^2}
    =
    c^2\nabla^2\bm{u}.
    }
\end{equation}

Thus the familiar wave equation can also be written for the displacement vector.

\section{Why pressure is usually preferred in acoustics}

Although displacement is physically meaningful in a fluid, pressure is often more convenient.

A three-dimensional displacement field is a vector:

\begin{equation}
    \bm{u}=
    \begin{pmatrix}
    u_x\\
    u_y\\
    u_z
    \end{pmatrix}.
\end{equation}

Pressure, on the other hand, is a scalar:

\begin{equation}
    p=p(\bm{x},t).
\end{equation}

Because an ideal acoustic fluid supports only compressional waves, one scalar variable is sufficient to describe the propagating mechanical mode.

This is why pressure is commonly used in acoustic modeling even though the actual material particles still undergo vector displacement.

\begin{tcolorbox}[title=Useful comparison]
\[
\boxed{
\begin{array}{ccl}
\text{Elastic solid}
&:&
\bm{u}\ \text{is natural because compression and shear are present},
\\[2mm]
\text{Acoustic fluid}
&:&
\bm{u}\ \text{is valid, but only its compressional part propagates},
\\[2mm]
&&
p\ \text{provides a convenient scalar description}.
\end{array}
}
\]
\end{tcolorbox}

\section{Connection between elasticity and acoustics}

For an isotropic elastic solid, the stress tensor is

\begin{equation}
    \bm{\sigma}
    =
    \lambda(\nabla\cdot\bm{u})\bm{I}
    +
    2\mu\bm{\epsilon},
\end{equation}

where $\lambda$ and $\mu$ are the Lam\'e parameters.

The parameter $\mu$ is the shear modulus. A fluid cannot sustain shear stress, so in the acoustic limit,

\begin{equation}
    \mu=0.
\end{equation}

The stress tensor therefore becomes purely isotropic:

\begin{equation}
    \bm{\sigma}
    =
    K(\nabla\cdot\bm{u})\bm{I},
\end{equation}

where $K$ is the bulk modulus.

Fluid stress is conventionally written as

\begin{equation}
    \bm{\sigma}=-p\bm{I}.
\end{equation}

Comparing the two expressions gives

\begin{equation}
    \boxed{
    p=-K\nabla\cdot\bm{u}.
    }
\end{equation}

This provides a useful conceptual bridge:

\begin{equation}
    \boxed{
    \text{acoustic wave propagation can be viewed as the compressional limit of elasticity.}
    }
\end{equation}

\section{Summary}

The main ideas are:

\begin{enumerate}
    \item $\bm{u}$ is the displacement of a material particle from its equilibrium position, not its absolute position.

    \item Strain depends on spatial derivatives of $\bm{u}$ because deformation is caused by relative displacement between neighboring particles.

    \item The particle velocity is
    \[
        \bm{v}=\frac{\partial\bm{u}}{\partial t}.
    \]

    \item Displacement can be used to describe waves in both solids and fluids.

    \item In a fluid, only volumetric deformation is relevant:
    \[
        \nabla\cdot\bm{u}.
    \]

    \item The acoustic pressure perturbation is
    \[
        p=-K\nabla\cdot\bm{u}.
    \]

    \item The displacement-form acoustic wave equation is
    \[
        \rho\frac{\partial^2\bm{u}}{\partial t^2}
        =
        \nabla\left(K\nabla\cdot\bm{u}\right).
    \]

    \item Pressure is usually preferred in acoustic modeling because it reduces the problem to a scalar field.
\end{enumerate}

\begin{tcolorbox}[title=Central physical interpretation]
The particles of a medium do not travel with the wave over large distances. They oscillate around their equilibrium positions.

The displacement field $\bm{u}$ describes those local perturbations, while the wave equation describes how those perturbations propagate through the medium.
\end{tcolorbox}

\end{document}
